% Generated draft: requires mathematical and source review.
\documentclass[11pt]{article}
\usepackage[margin=25mm]{geometry}
\usepackage{fontspec}
\setmainfont{Latin Modern Roman}
\usepackage{amsmath,amssymb,mathtools,unicode-math}
\setmathfont{Latin Modern Math}
\usepackage{enumitem,needspace,xurl,hyperref}
\setlength{\parindent}{0pt}
\setlength{\parskip}{5pt}
\setlength{\emergencystretch}{3em}
\newcommand{\N}{\mathbb{N}}
\newcommand{\Z}{\mathbb{Z}}
\newcommand{\Q}{\mathbb{Q}}
\newcommand{\R}{\mathbb{R}}
\newcommand{\C}{\mathbb{C}}
\newcommand{\F}{\mathbb{F}}
\newcommand{\A}{\mathbb{A}}
\newcommand{\PP}{\mathbb{P}}
\newcommand{\E}{\mathbb{E}}
\newcommand{\eps}{\varepsilon}
\newcommand{\dd}{\,\mathrm{d}}
\begin{document}
\section*{UnsolvedMath ID: 30002085}
\textbf{Source problem code:} catalogue entry
\subsection{Definitions and mathematical statement}
For this draft, a Ricci-flat ALE four-manifold means a connected, complete, smooth, noncompact Riemannian manifold $(M,g)$ of real dimension $4$, without boundary, with $\operatorname{Ric}_g=0$, satisfying the following conventional asymptotically locally Euclidean condition. There exist a compact set $K\subset M$, a number $R>0$, a finite subgroup $\Gamma\subset O(4)$ acting freely on $S^3$, and a diffeomorphism
\[
\Phi:(\R^4\setminus\overline{B_R(0)})/\Gamma\longrightarrow M\setminus K
\]
such that, after lifting $\Phi^*g$ to $\R^4\setminus\overline{B_R(0)}$, there is some $\tau>0$ for which
\[
\partial^\alpha\bigl((\Phi^*g)_{ij}-\delta_{ij}\bigr)
=O\bigl(r^{-\tau-|\alpha|}\bigr)
\qquad(r\to\infty)
\]
for every multi-index $\alpha$. Here $r$ is the Euclidean radius and $\delta$ is the Euclidean metric. Completeness, the end convention, the decay and differentiability requirements, and the allowed group at infinity are conventional choices made explicit here; the supplied problem does not specify them or explicitly impose orientability.

\medskip
Let $\mathcal K$ denote the family referred to in the source as Kronheimer gravitational instantons, and let $\mathcal Q_{\mathrm{source}}$ denote the family referred to there as their known finite quotients. A smooth finite Riemannian quotient $X/G$ means a quotient by a finite group $G$ acting freely by isometries, equipped with the induced metric. The source does not enumerate $\mathcal Q_{\mathrm{source}}$, specify the exact scope of $\mathcal K$, or say whether every smooth finite quotient is intended.

\textbf{Source question, with the unresolved family qualification retained.} Does there exist a Ricci-flat ALE four-manifold $(M,g)$ such that
\[
(M,g)\not\cong X
\qquad\text{for every }X\in\mathcal K\cup\mathcal Q_{\mathrm{source}},
\]
where $\cong$ denotes Riemannian isometry? This expresses the supplied existence question, but it does not define a fully enumerated exclusion class: the phrase ``known finite quotients'' requires clarification from the source before the target can be made completely precise. The treatment of flat examples and the intended equivalence relation also require confirmation.

\medskip
\textbf{Related question in the supplied dated curation.} Must every simply connected Ricci-flat ALE four-manifold be hyperkähler? Here hyperkähler means that $M$ admits three parallel, metric-compatible complex structures $I,J,L$ satisfying
\[
I^2=J^2=L^2=-\operatorname{id},\qquad IJ=L=-JI.
\]
The simply connected hypothesis belongs to this related formulation; it is not imposed in the main source question. The supplied evidence does not establish equivalence between these two formulations. The curation describes Kronheimer's classification as covering hyperkähler ALE four-manifolds and mentions results under additional hypotheses, without specifying those hypotheses. Those assertions are retained as supplied background, not as independently verified classification statements.
\subsection{Short English statement}
The source asks whether a complete Ricci-flat ALE four-manifold can exist outside Kronheimer's gravitational instantons and their known finite quotients. The excluded quotient family and the exact ALE conventions are not specified. The supplied curation also asks whether every simply connected example is hyperkähler; this related question is not asserted to be equivalent to the original one. The supplied open status requires independent review.
\subsection{Sources}
\begin{itemize}
\item Local UnsolvedMath registry, permanent ID \texttt{30002085}, source problem code \texttt{catalogue entry}; collection: Oberwolfach Reports Open Problems. Repository registry references: \texttt{data/unsolvedmath/problems.json} and the supplementary \texttt{catalog/problems.json}. This draft uses only the supplied complete record.
\item Supplied report citation: \emph{Geometrie (2013)}, Oberwolfach Reports, DOI \texttt{10.4171/owr/2012/27}, pp.~1639--1686. \url{https://doi.org/10.4171/owr/2012/27}. The record gives proposed year 2012 and does not identify an individual proposer. The report-title year and proposed year are preserved without reconciliation.
\item Supplied literature source: Li, \emph{On 4-dimensional Ricci-flat ALE manifolds}. \url{https://arxiv.org/abs/2304.01609}.
\item Supplied literature source: Li--Sun, \emph{Geometry of gravitational instantons}. \url{https://arxiv.org/abs/2606.10443}.
\item Supplied literature source, labeled ``original Oberwolfach report'': \url{https://doi.org/10.4171/owr/2012/27}.
\item Literature assessment dated \texttt{2026-08-21}: the supplied curation labels the problem open, describes Kronheimer's classification as covering hyperkähler ALE four-manifolds, mentions results under additional hypotheses, and says that the simply connected Ricci-flat question remains open. This is supplied dated curation, not an independently verified literature assessment. The record itself requests independent verification of time-sensitive status and cited claims.
\item The supplied statement status is \texttt{corrected\_verified}; its verification note attributes the existence question to the official report. This label and attribution are source metadata, not verification performed for this draft.
\end{itemize}
\end{document}
